\documentclass[border=2pt]{standalone}
\usepackage{amsmath, amsthm, amsfonts}
\usepackage{tikz-cd, kotex}
\usepackage{ytableau}
\usepackage{quiver}
\usepackage{Operators}
\usepackage{palette}
\usepackage{diagram-style}
\usetikzlibrary{arrows.meta, positioning, decorations.pathreplacing, decorations.markings, intersections}
\begin{document}

%%FIG diagram (Cluster_Monomial_Basis-1.svg)
\begin{tikzpicture}[scale=0.55]
  \draw[black!50, very thick, dashed] (0,0) rectangle (3,-3);
  \fill[black!20] (0,0) rectangle (3,-2);
  \draw (0,0) grid (3,-2);
  \fill[accent2!25] (0,-1) rectangle (1,-2);
  \fill[accent1!30] (1,0) rectangle (2,-1);
  \draw[very thick] (1,-1) rectangle (2,-2);
  \draw (0,0) grid (3,-2);
  \node at (1.5,-1.5) {\scriptsize $(i,j)$};
  \node[accent2!60!black] at (0.5,-2.6) {\scriptsize row 왼쪽};
  \node[accent1!70!black] at (2.7,0.5) {\scriptsize col 위쪽};
  \node at (1.5,-3.7) {\scriptsize $\lambda=(n\!-\!k)^a$의 cohook};
\end{tikzpicture}

\end{document}
